% lectura-previa-cap01.tex
% Lectura previa a la Sesión 1 — Unidad I: Comportamiento del consumidor
% Laboratorio: Gestión de datos, comportamiento y perfiles del consumidor digital
% Lic. en Mercadotecnia Digital · ESCA Unidad Santo Tomás · IPN
% Prof. D.Sc. Barsekh-Onji Aboud
\documentclass[12pt,letterpaper]{article}

\usepackage[utf8]{inputenc}
\usepackage[spanish,es-noshorthands]{babel}
\usepackage[T1]{fontenc}
\usepackage{geometry}
\geometry{top=2.5cm, bottom=2.5cm, left=2.8cm, right=2.8cm}
\usepackage{csquotes}
\usepackage{parskip}
\usepackage{enumitem}
\usepackage{xcolor}
\usepackage{hyperref}
\usepackage{fancyhdr}
\usepackage{tcolorbox}
\tcbuselibrary{skins,breakable}

\definecolor{guinda}{HTML}{6F1D46}

\hypersetup{
  colorlinks=true,
  linkcolor=guinda,
  urlcolor=guinda,
  pdftitle={Lectura previa -- Antes de empezar: quien es el consumidor digital},
  pdfauthor={Dr. Aboud Barsekh-Onji}
}

\setlength{\headheight}{14pt}
\pagestyle{fancy}
\fancyhf{}
\lhead{\small Laboratorio: Gestión de Datos, Comportamiento y Perfiles del Consumidor Digital}
\lfoot{\small IPN --- ESCA Unidad Santo Tomás}
\rfoot{\small Página \thepage}

\setlength{\parskip}{0.55\baselineskip}
\setlength{\parindent}{0pt}

\newtcolorbox{cajaInstrucciones}{
  colback=guinda!6, colframe=guinda, boxrule=0.6pt, arc=1mm,
  left=6pt, right=6pt, top=5pt, bottom=5pt, breakable
}

\begin{document}
\thispagestyle{fancy}

\begin{center}
  {\LARGE \textbf{Antes de empezar: \\[2pt] ¿quién es, en realidad, el consumidor digital?}}\\[6pt]
  {\large Lectura previa --- Unidad I, sesión 1}\\[4pt]
  {\small Lic. en Mercadotecnia Digital $\cdot$ ESCA Unidad Santo Tomás $\cdot$ IPN}\\[4pt]
  {\small Dr. Aboud Barsekh-Onji}
\end{center}

\vspace{0.3cm}

\begin{cajaInstrucciones}
\textbf{Antes de la clase.} Lean este texto completo (15--20 minutos) y traigan
por escrito, aunque sea en una frase cada una, sus respuestas a las
\textbf{preguntas de reflexión} del final. No hay respuesta \enquote{correcta}
todavía: el objetivo es que lleguen a la sesión 1 con una postura propia que
podamos poner a prueba juntos frente al tema \textbf{1.1, Comportamiento del
consumidor}.
\end{cajaInstrucciones}

\vspace{0.2cm}

Abran su teléfono y reconstruyan, mentalmente, su última compra en línea.
Probablemente no decidieron de golpe: compararon precios en dos o tres
aplicaciones, leyeron un par de reseñas, quizás vieron que un conocido ya
tenía el mismo producto, y solo entonces pagaron. Ese recorrido --- comparar,
verificar, imitar, decidir --- se siente completamente normal. Vale la pena
detenerse ahí, porque ese recorrido \enquote{normal} es exactamente el objeto de
estudio de esta unidad de aprendizaje, y la primera sesión empieza preguntando
qué tan nuevo es en realidad.

\textbf{Una curva, o una colección de libros.} Existen dos maneras de contar
esta historia. La primera es una curva medible: en México, el porcentaje de
la población de seis años o más que usa internet pasó de 57.4\,\% en 2015 a
86.1\,\% en 2025 (INEGI, ENDUTIH). Es un cambio real, gradual, documentado
año con año. La segunda manera de contar la misma historia es editorial: una
serie de libros de mercadotecnia ha ido nombrando \enquote{versiones} sucesivas
del consumidor --- 1.0, 2.0, 3.0, y así hasta la más reciente ---, cada una
publicada justo cuando cierta tecnología se pone de moda. Ambas historias
existen y ambas son útiles, pero no son lo mismo: una se puede verificar con
una fuente estadística pública; la otra es, ante todo, un producto editorial
exitoso. Distinguir cuál de las dos está usando alguien --- un profesor, una
agencia de marketing, un video que les recomendó el algoritmo --- cuando habla
del \enquote{nuevo consumidor} es la primera habilidad crítica que esta unidad
les pide desarrollar.

\textbf{Persona y dato, al mismo tiempo.} Hay algo más de fondo en ese
recorrido de compra que describieron al inicio: cada paso --- cada búsqueda,
cada reseña leída, cada segundo que pasaron viendo una foto de producto ---
queda registrado en algún sistema. Ustedes siguen siendo, por supuesto, una
persona con una historia, un estado de ánimo y una razón particular para
comprar. Pero también son, simultáneamente, una fila de datos estructurados
en la base de una plataforma. Ninguna de las dos lecturas basta por sí sola:
una empresa que solo ve el dato pierde el contexto humano de la decisión;
una que solo ve a la persona no puede operar a la escala de un mercado
digital de millones de usuarios. Sostener esa tensión sin resolverla de
forma simplista es, en buena medida, de lo que trata esta unidad.

\textbf{\enquote{Inteligente} no es lo mismo que \enquote{sostenible}.} El discurso de
mercadotecnia digital tiende a mezclar dos ideas que conviene separar desde
ahora. Comprar de forma \emph{inteligente} significa decidir con información:
comparar, leer especificaciones, verificar antes de pagar. Comprar de forma
\emph{sostenible} significa que esa decisión tiene un menor costo ambiental y
social, sin importar qué tan informada haya estado. Las dos cosas pueden
coincidir --- y a veces coinciden ---, pero no son automáticamente lo mismo:
se puede comparar cinco aplicaciones para encontrar el envío más barato sin
que ninguna de las cinco opciones considere, siquiera, la huella del
transporte. Noten que esta distinción no es un tecnicismo académico: es
exactamente el tipo de confusión que un profesional de mercadotecnia digital
tiene que saber señalar cuando alguien la usa, con buena o mala intención,
para vender algo como \enquote{verde} sin evidencia.

\textbf{¿Y para qué sirve, entonces, analizar todo esto?} Una organización que
analiza sistemáticamente el comportamiento de su consumidor digital puede
segmentar con más precisión, personalizar su oferta y detectar en qué punto
del recorrido de compra la gente abandona. Eso está razonablemente bien
sustentado. Lo que está mucho menos sustentado --- y por eso conviene
sospechar de entrada --- es la promesa, muy común en el discurso comercial de
analítica de datos, de que esos mismos datos permiten \emph{anticipar}
tendencias futuras con confianza. Un patrón detectado en datos históricos no
garantiza que se vaya a repetir. Aprender a distinguir la ventaja de datos
que sí está probada de la que todavía es, sobre todo, un argumento de venta
de un proveedor de tecnología, es una de las competencias centrales que esta
unidad busca formar en ustedes.

Con estas tres ideas en la cabeza --- curva contra narrativa, persona y dato
a la vez, inteligente no es igual a sostenible --- están listos para entrar,
en la sesión 1, al tema \textbf{1.1, Comportamiento del consumidor}, donde
vamos a poner cada una de ellas a prueba con evidencia.

\vspace{0.4cm}
\hrule
\vspace{0.4cm}

\section*{Preguntas de reflexión}

Respondan brevemente, en sus propias palabras, antes de la sesión 1.

\begin{enumerate}[leftmargin=1.4em, itemsep=6pt]
  \item Piensen en su compra digital más reciente. ¿Qué parte de su decisión
  fue realmente independiente y qué parte estuvo influida por ver que otras
  personas ya habían comprado o recomendado lo mismo?
  \item ¿Habían escuchado hablar del \enquote{consumidor 4.0} o \enquote{5.0} antes
  de esta lectura? ¿De dónde les llegó esa idea, y alguna vez vieron el dato
  o el estudio detrás de la etiqueta, o solo el nombre?
  \item Den un ejemplo propio de una compra suya que haya sido \emph{informada}
  (comparada, investigada) pero probablemente no \emph{sostenible}. ¿Pueden
  pensar también en el caso contrario?
  \item Si en unos meses fueran responsables de mercadotecnia digital de una
  empresa mexicana pequeña, ¿cuál de estas ventajas de analizar al
  consumidor priorizarían primero: segmentar mejor, personalizar la oferta,
  reducir el abandono en el proceso de compra, o anticipar tendencias? ¿Por
  qué, y qué riesgo tiene esa elección?
  \item Si alguien les presentara una gráfica y dijera \enquote{esta tendencia va
  a seguir creciendo así}, ¿qué le preguntarían antes de creerle?
\end{enumerate}

\vspace{0.3cm}
\begin{center}
\small\textit{Guarden sus respuestas: seremos parte de la discusión que abre la siguiente sesión}
\end{center}

\end{document}
